\documentclass[11pt,letterpaper]{article}

\usepackage[spanish,es-noshorthands]{babel}
\usepackage{fontspec}
\setmainfont{Source Serif Pro}
\usepackage{microtype}
\usepackage{geometry}
\usepackage{booktabs}
\usepackage{tabularx}
\usepackage{enumitem}
\usepackage{xcolor}
\usepackage{graphicx}
\usepackage{float}
\usepackage{csquotes}
\usepackage{amsmath}
\usepackage{siunitx}
\usepackage{tikz}
\usetikzlibrary{arrows.meta,positioning,shapes.geometric,fit,backgrounds,decorations.pathreplacing}
\usepackage{xurl}
\usepackage[hidelinks]{hyperref}
\usepackage[backend=biber,style=apa,sorting=nyt]{biblatex}
\DeclareLanguageMapping{spanish}{spanish-apa}

\geometry{margin=2.3cm}
% Los rótulos de las figuras usan nodos de ancho fijo con texto alineado, y la
% bibliografía compone en modo sloppy: ambos generan avisos de línea holgada
% que son cosméticos. Se silencian para que `make check` siga siendo sensible
% a los desbordes reales.
\hbadness=10000
\setlength{\emergencystretch}{1em}
\setlength{\parindent}{0pt}
\setlength{\parskip}{0.42em}
\setlist{nosep}
\renewcommand{\topfraction}{0.92}
\renewcommand{\bottomfraction}{0.85}
\renewcommand{\textfraction}{0.06}
\renewcommand{\floatpagefraction}{0.72}
\setcounter{topnumber}{2}

\addbibresource{referencias.bib}
\AtBeginBibliography{\sloppy\small}
\setlength{\bibitemsep}{0.28em}
\setlength{\bibhang}{1.4em}

\sisetup{
  output-decimal-marker = {,},
  group-separator = {.},
  group-minimum-digits = 4,
  detect-all
}

% Colores: el azul del silicio propio, el gris del x86 al que sustituye, el
% verde de la mejora medida, el ámbar de la respuesta de los competidores y el
% rojo de lo que la adopción costó.
\definecolor{cApple}{RGB}{36,102,148}
\definecolor{cX86}{RGB}{104,110,126}
\definecolor{cVerde}{RGB}{56,120,86}
\definecolor{cAmbar}{RGB}{166,92,34}
\definecolor{cAviso}{RGB}{176,42,42}
\definecolor{cGris}{RGB}{132,124,112}

% --- Panel A: rendimiento multinúcleo frente a potencia de paquete ----------
\newcommand{\px}[1]{(#1)*0.205}
\newcommand{\py}[1]{(#1)*0.19}   % el valor se pasa en miles de puntos
% #1 x (vatios), #2 y (puntos), #3 color, #4 posición del rótulo, #5 rótulo
\newcommand{\punto}[5]{%
  \fill[#3] ({\px{#1}},{\py{#2}}) circle (2.5pt);%
  \node[#4, font=\scriptsize, text=#3, inner sep=3.2pt, align=center]
       at ({\px{#1}},{\py{#2}}) {#5};%
}

% --- Panel B: trayectoria temporal de la figura de mérito -------------------
\newcommand{\tx}[1]{((#1)-2020)*0.98}
\newcommand{\ty}[1]{(#1)*0.19}   % el valor se pasa en miles de puntos
% #1 año, #2 puntos, #3 color, #4 posición, #5 rótulo
\newcommand{\hito}[5]{%
  \fill[#3] ({\tx{#1}},{\ty{#2}}) circle (2.3pt);%
  \node[#4, font=\scriptsize, text=#3, inner sep=3pt] at
       ({\tx{#1}},{\ty{#2}}) {#5};%
}

\title{\vspace{-1.2cm}\textbf{El paquete lo cambió todo}\\
\large Impacto neto de la adopción de Apple Silicon en el Mac, del M1 (2020)
al M5 (2025), sobre el producto, sus consumidores, sus competidores y el
mercado}
\author{Carlos Luis Rojas Aragonés\\
\small Gestión del Desarrollo Tecnológico: Estrategias y Análisis de Portfolio}
\date{\small 16 de septiembre de 2026}

\begin{document}

\maketitle
\vspace{-0.8em}

\section{El caso, en una ficha}

El ejemplo que traigo es la sustitución de los procesadores x86-64 de Intel
por un sistema en chip propio de arquitectura Arm en toda la línea Mac,
anunciada en junio de 2020 y aún en curso cinco generaciones después
\autocite{apple2020transicion}. Es un buen caso precisamente por lo que señala
el planteamiento del módulo: no fue un cambio de proveedor, fue un cambio de
arquitectura, y el ejercicio consiste en separar cuánto del resultado viene
del silicio y cuánto de haber movido una interfaz de sitio.

\subsection*{Lo que encontré, dicho por adelantado}

\begin{quote}
\textbf{Lo que se adoptó no fue un procesador más rápido: fue la desaparición
de la interfaz de memoria del sistema.} El rendimiento por vatio es la
consecuencia, no la causa. Frente al Mac Intel al que sustituyó, el M1
multiplicó por 2,0 la puntuación multinúcleo y por 2,7 el rendimiento por
vatio, y la línea siguió creciendo al 16,5 por ciento anual hasta el M5. Pero
el dato que mejor mide el impacto neto no está en el Mac: está en que Intel
integró la memoria en el paquete cuatro años después, Qualcomm entró en el PC
con los propios arquitectos que habían diseñado esos núcleos y Microsoft
redefinió su gama en torno a una métrica que en 2020 no existía. \textbf{Y el
coste fue real}: el 90 por ciento de los elementos del sistema hubo que
tocarlos, y el comprador perdió la memoria ampliable, la GPU externa y el
arranque nativo de Windows.
\end{quote}

\begin{table}[H]
\small
\caption{La adopción, en corto.}
\label{cuadro:ficha}
\begin{tabularx}{\textwidth}{@{}p{4.3cm}X@{}}
\toprule
\textbf{Dimensión} & \textbf{Dato} \\
\midrule
Producto de referencia
& Mac con procesador Intel, GPU discreta AMD en la gama alta, memoria fuera
  del paquete y coprocesador T2 \\
\addlinespace[2pt]
Tecnología adoptada
& Sistema en chip propio (AArch64) con CPU heterogénea, GPU, Neural Engine,
  codificadores de vídeo y memoria LPDDR en el mismo paquete \\
\addlinespace[2pt]
Calendario
& Anuncio en junio de 2020; primeros equipos con M1 el 10 de noviembre de
  2020; transición cerrada en junio de 2023, 31 meses después
  \autocite{apple2020m1,apple2023macpro} \\
\addlinespace[2pt]
Alcance del análisis
& Cinco generaciones, de M1 (5 nm, 2020) a M5 (tercera generación de 3 nm,
  octubre de 2025) \autocite{apple2025m5} \\
\addlinespace[2pt]
Figuras de mérito
& Principal: rendimiento multinúcleo sostenido por vatio de paquete.
  Secundarias: ancho de banda de memoria, autonomía real y coste por punto de
  rendimiento para el comprador \\
\bottomrule
\end{tabularx}
\end{table}

\subsection*{Nota sobre las cifras}

Separo tres clases de número, porque mezclarlas haría parecer el análisis más
sólido de lo que es. \textbf{Publicadas por el fabricante}: fechas, proceso de
fabricación, transistores, ancho de banda de memoria, autonomía declarada,
precios de lista y ventas netas del segmento Mac de los informes anuales
\autocite{apple10k}. \textbf{Medidas por terceros}: las puntuaciones de
Geekbench 6 son medianas de resultados públicos, con una dispersión del orden
del 5 por ciento, y las uso solo para comparar órdenes de magnitud
\autocite{geekbench2026}. \textbf{Estimaciones propias}, marcadas donde
aparecen: las potencias de paquete sostenidas, el recuento de elementos de la
descomposición del sistema, el coste por punto de rendimiento y las tasas de
crecimiento derivadas.

\section{Qué se sustituyó exactamente}

El caso del sistema de impresión digital que vimos en el módulo mide el
esfuerzo de adopción sobre la matriz de estructura de diseño y llega al 8,5
por ciento de cambio respecto del sistema de referencia. Aplicando la misma
lógica de matriz delta a una descomposición del Mac en unos 32 elementos entre
componentes, interfaces y capas de software, el resultado de esta adopción es
del orden del 90 por ciento (Cuadro~\ref{cuadro:delta}). Es la diferencia
entre una mejora de componente y lo que \textcite{henderson1990} llaman
innovación arquitectónica: los componentes siguen haciendo lo mismo, pero la
forma en que se conectan cambia por completo.

\begin{table}[H]
\small
\caption{Matriz delta de la adopción, en la línea del paso 3 del caso del
módulo. El recuento de elementos es elaboración propia sobre una
descomposición de 32 elementos y sirve para dimensionar el esfuerzo, no como
inventario exhaustivo.}
\label{cuadro:delta}
\begin{tabularx}{\textwidth}{@{}p{2.55cm}p{1.15cm}X@{}}
\toprule
\textbf{Categoría} & \textbf{N.\textsuperscript{o}} & \textbf{Elementos} \\
\midrule
Añadidos & 9
& Sistema en chip unificado; memoria LPDDR dentro del paquete; Neural Engine;
  codificadores de vídeo dedicados; Secure Enclave integrado; Rosetta 2;
  binarios universales; cadena de arranque con política local; virtualización
  nativa de sistemas Arm \\
\addlinespace[2pt]
Eliminados & 6
& Procesador Intel; chipset acompañante; coprocesador T2; GPU discreta AMD;
  módulos de memoria fuera del paquete; Boot Camp y soporte de GPU externa \\
\addlinespace[2pt]
Modificados & 14
& Placa base y alimentación; sistema térmico (el MacBook Air pierde el
  ventilador); Thunderbolt pasa a integrarse; pila gráfica Metal, que deja de
  copiar datos entre memorias; controladores de terceros, de extensiones de
  núcleo a DriverKit; cadena de compilación; firma y notarización; macOS 11 \\
\addlinespace[2pt]
Sin cambio & 3
& Pantalla, teclado y trackpad, es decir, casi todo lo que el usuario toca \\
\bottomrule
\end{tabularx}
\end{table}

La interfaz que desaparece es la clave técnica de todo lo demás. En el sistema
de referencia, la CPU hablaba con la DRAM por un bus externo y con la GPU por
PCI Express, y cada copia de datos entre ambas costaba tiempo y energía. Al
poner la memoria en el paquete y hacerla visible para CPU, GPU y Neural Engine
a la vez, esa copia deja de existir. Por eso el ancho de banda crece de unos
60 GB/s en el portátil Intel de referencia a 68 GB/s en el M1 y a 153 GB/s en
el M5, y en las variantes Max y Ultra a 400 y 800 GB/s, cifras que un portátil
de memoria modular no puede alcanzar sin pagar un coste energético que no cabe
en el chasis.

\section{Las figuras de mérito}

\begin{table}[b]
\small
\caption{Figuras de mérito del silicio. La referencia Intel es el procesador
del MacBook Pro de 13 pulgadas de 2020, el equipo al que sustituyó
directamente el primer M1. Las potencias de paquete y la fila de rendimiento por vatio
son estimaciones propias. La última fila compara MacBook Air con MacBook Air,
y su columna de 2025 corresponde al modelo con M4, el más reciente de esa
gama con autonomía publicada.}
\label{cuadro:fom}
\begin{tabularx}{\textwidth}{@{}Xrrrr@{}}
\toprule
\textbf{Figura de mérito} & \textbf{Intel 2020} & \textbf{M1 2020} &
\textbf{M5 2025} & \textbf{M5/Intel} \\
\midrule
Geekbench 6, un núcleo (puntos) & 1.550 & 2.350 & 4.130 & 2,7$\times$ \\
Geekbench 6, multinúcleo (puntos) & 4.300 & 8.400 & 17.800 & 4,1$\times$ \\
Ancho de banda de memoria (GB/s) & 60 & 68 & 153 & 2,6$\times$ \\
Proceso de fabricación (nm) & 10 & 5 & 3 & \\
Potencia de paquete sostenida (W) & 28 & 20 & 25 & 0,9$\times$ \\
\addlinespace[2pt]
\textbf{Multinúcleo por vatio (puntos/W)} & \textbf{154} & \textbf{420} &
\textbf{712} & \textbf{4,6}$\times$ \\
Autonomía de vídeo del MacBook Air (h) & 12 & 18 & 18 & 1,5$\times$ \\
\bottomrule
\end{tabularx}
\end{table}

La figura de mérito que ordena el caso es el rendimiento multinúcleo por vatio
de paquete, porque es la que fija a la vez la autonomía, el volumen del chasis
y el ruido, que son los tres atributos por los que el cliente paga en un
portátil. El Cuadro~\ref{cuadro:fom} la sitúa junto a las demás. El salto
inicial es de 154 a 420 puntos por vatio, un factor de 2,7, y el acumulado
hasta el M5 es de 4,6.

Sobre la trayectoria posterior, que es la parte del ejemplo que más interesa
al usuario del producto, la tasa de crecimiento compuesta de la figura de
mérito multinúcleo entre el M1 (noviembre de 2020) y el M5 (octubre de 2025),
es decir 4,92 años, es

\begin{equation}
g \;=\; \left(\frac{17.800}{8.400}\right)^{1/4,92} - 1 \;=\; 0{,}165
\qquad\Longrightarrow\qquad
t_{2} \;=\; \frac{\ln 2}{\ln(1+g)} \;=\; 4{,}5 \text{ años}
\label{eq:cagr}
\end{equation}

un 16,5 por ciento anual y una duplicación cada cuatro años y medio. Conviene
ser honesto con lo que esto significa: es una tasa buena y sostenida, pero no
es extraordinaria en términos históricos, y la parte del salto que impresionó
en 2020 no se repitió. El cambio de arquitectura se cobra una sola vez.

\begin{figure}[!htb]
\centering
\begin{tikzpicture}

% ----- Panel A -----
\begin{scope}
  \node[anchor=west, font=\small\bfseries] at (-1.5,4.35)
       {A. Dónde está cada silicio};
  \node[anchor=west, font=\scriptsize, text=cGris] at (-1.5,3.95)
       {Multinúcleo Geekbench 6 frente a potencia de paquete estimada};
  % ejes
  \draw[cGris, thin] (0,0) -- (6.9,0);
  \draw[cGris, thin] (0,0) -- (0,3.75);
  \foreach \w in {10,20,30}
    \draw[cGris, thin] ({\px{\w}},0) -- ({\px{\w}},-0.08)
      node[below, font=\scriptsize, text=cGris] {\w\,W};
  \foreach \p/\r in {5/5.000,10/10.000,15/15.000}
    \draw[cGris, thin] (0,{\py{\p}}) -- (-0.08,{\py{\p}})
      node[left, font=\scriptsize, text=cGris] {\r};
  % rectas de isoeficiencia
  \draw[cX86, dashed, thin] (0,0) -- (6.9,{6.9*0.1427});
  \node[font=\scriptsize, text=cX86, rotate=8.1] at (5.85,0.71)
       {154 puntos/W};
  \draw[cApple, dashed, thin] (0,0) -- (5.3,{5.3*0.660});
  \node[font=\scriptsize, text=cApple, rotate=33.4] at (2.75,2.02)
       {712 puntos/W};
  % puntos
  \punto{28}{4.3}{cX86}{above left}{Intel 2020\\\tiny i7-1068NG7}
  \punto{30}{10.7}{cX86}{right}{Intel 2024}
  \punto{28}{15.0}{cAmbar}{above}{AMD 2024}
  \punto{23}{14.0}{cAmbar}{below right}{Qualcomm 2024}
  \punto{20}{8.4}{cApple}{left}{M1}
  \punto{22}{14.8}{cApple}{above left}{M4}
  \punto{25}{17.8}{cApple}{above}{M5}
\end{scope}

% ----- Panel B -----
\begin{scope}[xshift=9.15cm]
  \node[anchor=west, font=\small\bfseries] at (-1.5,4.35)
       {B. La trayectoria, generación a generación};
  \node[anchor=west, font=\scriptsize, text=cGris] at (-1.5,3.95)
       {Multinúcleo Geekbench 6, 2020 a 2025};
  \draw[cGris, thin] (0,0) -- (6.0,0);
  \draw[cGris, thin] (0,0) -- (0,3.75);
  \foreach \a in {2020,2022,2024}
    \draw[cGris, thin] ({\tx{\a}},0) -- ({\tx{\a}},-0.08)
      node[below, font=\scriptsize, text=cGris] {\a};
  \foreach \p/\r in {5/5.000,10/10.000,15/15.000}
    \draw[cGris, thin] (0,{\py{\p}}) -- (-0.08,{\py{\p}})
      node[left, font=\scriptsize, text=cGris] {\r};
  % serie Apple
  \draw[cApple, thick]
    ({\tx{2020.85}},{\ty{8.4}}) -- ({\tx{2022.45}},{\ty{9.7}}) --
    ({\tx{2023.8}},{\ty{11.7}}) -- ({\tx{2024.8}},{\ty{14.8}}) --
    ({\tx{2025.8}},{\ty{17.8}});
  \hito{2020.85}{8.4}{cApple}{above left}{M1}
  \hito{2022.45}{9.7}{cApple}{below right}{M2}
  \hito{2023.8}{11.7}{cApple}{below right}{M3}
  \hito{2024.8}{14.8}{cApple}{below right}{M4}
  \hito{2025.8}{17.8}{cApple}{above left}{M5}
  % serie x86 de referencia
  \draw[cX86, thick, dashed]
    ({\tx{2020.2}},{\ty{4.3}}) -- ({\tx{2024.7}},{\ty{10.7}});
  \hito{2020.2}{4.3}{cX86}{below right}{x86 de 2020}
  \hito{2024.7}{10.7}{cX86}{below right}{x86 de 2024}
  % anotación de la tasa
  \node[font=\scriptsize, text=cVerde, align=left, anchor=west]
       at (0.35,3.30) {16,5\,\% anual, duplicación cada 4,5 años};
  \node[font=\scriptsize, text=cX86, align=left, anchor=west]
       at (2.2,1.05) {23\,\% anual desde la mitad};
\end{scope}

\end{tikzpicture}
\caption{Panel A: la posición de cada silicio en el plano rendimiento
multinúcleo frente a potencia. Las rectas punteadas son isoeficiencias; cuanto
más inclinada, más rendimiento por vatio. Panel B: la misma figura de mérito a
lo largo del tiempo. La línea x86 crece más deprisa en términos relativos,
pero desde la mitad del nivel, y lo hace después de adoptar la memoria en el
paquete. Las potencias del Panel A son estimaciones propias con una
incertidumbre del orden de $\pm$20 por ciento; las puntuaciones son medianas
públicas \autocite{geekbench2026}.}
\label{fig:fom}
\end{figure}

\section{El impacto sobre el producto y sobre el consumidor}

\textbf{Sobre el producto.} El MacBook Air pasó de 12 a 18 horas de
reproducción de vídeo y perdió el ventilador, lo que elimina una pieza móvil,
una entrada de polvo y el modo de fallo más común del chasis. El MacBook Pro
de 13 pulgadas pasó de 10 a 20 horas de navegación. En la gama alta, la
memoria unificada de 400 y 800 GB/s permitió ejecutar en un portátil cargas de
trabajo (vídeo de alta resolución y, después, modelos de lenguaje locales) que
antes exigían una estación de sobremesa con GPU discreta.

\textbf{Sobre el consumidor, en beneficio.} El precio de lista no subió: el
MacBook Air costaba 999 dólares con Intel en 2020, 999 con M1 ese mismo
noviembre y 999 en 2025, ya con 16 GB de memoria unificada de serie
\autocite{apple2025air}. Manteniendo el precio y usando la puntuación
multinúcleo como medida del producto entregado, el coste por punto de
rendimiento del equipo de entrada cae de unos 0,42 dólares en el Air Intel de
2020 a 0,12 con el M1 y a 0,068 con el M4 de 2025: una reducción del 84 por
ciento en cinco años, con el precio nominal congelado.

\textbf{Sobre el consumidor, en coste.} Sería deshonesto presentar solo el
lado bueno. La adopción eliminó la memoria y el almacenamiento ampliables, la
GPU externa y el arranque nativo de Windows, de modo que el comprador cambió
modularidad por integración y quedó atado a la configuración que eligió el
día de la compra, en un producto cuyo incremento de memoria se factura muy por
encima del precio del componente. Durante la transición, el software x86 se
ejecutó traducido por Rosetta 2 con una penalización típica del 20 al 30 por
ciento, y hubo categorías enteras (virtualización de x86, algunas cadenas de
herramientas científicas, controladores de periféricos especializados) que
tardaron dos años o más en estar disponibles, cuando lo estuvieron.

\section{El impacto sobre los competidores y sobre el mercado}

\textbf{Intel} perdió una cuenta del orden de 20 millones de unidades anuales,
pero el efecto relevante no es esa pérdida sino el cambio de eje competitivo.
En 2024, con la plataforma Core Ultra 200V, Intel integró la memoria LPDDR5X
en el propio paquete del procesador \autocite{intel2024lunarlake}, que es
exactamente la decisión arquitectónica del M1 cuatro años después, y la
abandonó en la generación siguiente por su coste, lo que ilustra que la
integración no es gratis para quien vende silicio a terceros y no controla el
producto final. \textbf{AMD} respondió ensanchando el bus de memoria de sus
procesadores con gráficos integrados hasta niveles antes reservados a las GPU
discretas. \textbf{Qualcomm} entró en el mercado del PC en 2024 con núcleos
diseñados por el equipo de Nuvia, fundado por antiguos arquitectos de los
núcleos de Apple \autocite{qualcomm2024}. Y \textbf{Microsoft} redefinió su
gama alrededor de los Copilot+ PC exigiendo una unidad neuronal de al menos 40
billones de operaciones por segundo \autocite{microsoft2024}, es decir, una
figura de mérito que en 2020 no formaba parte de la conversación del PC y que
el M1 ya traía de fábrica.

\textbf{Sobre el mercado}, las ventas netas del segmento Mac pasaron de 28.622
millones de dólares en el ejercicio 2020 a 35.190 en 2021 y 40.177 en 2022, un
crecimiento del 40 por ciento en dos ejercicios, con una cuota mundial de
unidades que subió de en torno al 7 al entorno del 10 por ciento
\autocite{apple10k,idc2026}. El dato que evita la lectura triunfalista es el
del ejercicio 2023: 29.357 millones, una caída del 27 por ciento. La adopción
no inmunizó al producto frente al ciclo del mercado de PC ni frente a la
resaca de la compra masiva de la pandemia. Lo que sí hizo fue mover a Apple
del segmento en el que competía por gigahercios a uno en el que compite por
autonomía y rendimiento por vatio, donde su integración vertical es una
ventaja estructural y no una preferencia de diseño.

\section{El impacto neto}

Reuniendo las cuatro dimensiones que pide el planteamiento, el balance es el
siguiente. \textbf{Producto}: 2,0 veces el rendimiento multinúcleo y 2,7 veces
el rendimiento por vatio respecto del equipo al que sustituyó, con 4,6 veces
acumuladas hasta el M5 y un crecimiento posterior del 16,5 por ciento anual.
\textbf{Consumidor}: el mismo precio nominal durante cinco años, un 84 por
ciento menos de coste por punto de rendimiento y un 50 por ciento más de
autonomía, a cambio de perder modularidad, compatibilidad nativa con Windows y
la posibilidad de ampliar el equipo después de comprarlo. \textbf{Competidores}:
la pérdida de una cuenta grande para Intel y, sobre todo, la obligación de
copiar la decisión arquitectónica y de entrar en una carrera de eficiencia y
de capacidad neuronal que no habían elegido. \textbf{Mercado}: un 40 por
ciento más de ingresos del segmento en dos ejercicios, tres puntos de cuota y
un cambio en la figura de mérito con la que el sector entero se compara.

Y el contraste con el caso del módulo es, para mí, la lección del ejercicio.
El sistema de impresión digital cambió el 8,5 por ciento de su matriz para
ganar un 1,3 por ciento de valor percibido. Esta adopción cambió en torno al
90 por ciento de la suya para multiplicar su figura de mérito por 2,7 de
golpe. Las dos decisiones pueden ser correctas: lo que no es comparable es el
riesgo. La segunda solo era defendible porque el esfuerzo de adopción estaba
casi todo amortizado antes de empezar, en los diez años y los más de dos mil
millones de unidades de la serie A del iPhone y del iPad. Quien no tenga esa
línea base previa no puede leer este caso como una receta.

\section{Limitaciones}

\begin{enumerate}
\item \textbf{Las potencias de paquete son estimaciones propias}, no medidas
  de laboratorio, y la figura de mérito central del Cuadro~\ref{cuadro:fom}
  depende directamente de ellas. Un error del 20 por ciento en el denominador
  mueve el factor de 4,6 a un rango de 3,8 a 5,8.
\item \textbf{Geekbench 6 no mide el trabajo del usuario.} Es una medida
  sintética, comparable entre arquitecturas con reservas, y deja fuera la
  penalización de la emulación, que durante los dos primeros años fue la
  experiencia real de buena parte del software.
\item \textbf{El coste por punto de rendimiento supone que el rendimiento es
  lo que el comprador paga}, y no lo es: paga pantalla, chasis, sistema
  operativo y servicio. La cifra sirve para comparar un producto consigo mismo
  a lo largo del tiempo, no para comparar productos distintos.
\item \textbf{El recuento de elementos de la matriz delta es propio.} Con una
  descomposición más fina o más gruesa el porcentaje se mueve; lo que no se
  mueve es el orden de magnitud frente al 8,5 por ciento del caso del módulo.
\item \textbf{No puedo separar el efecto del cambio de arquitectura del
  efecto del proceso de fabricación.} Entre 2020 y 2025 hubo tres saltos de
  nodo, de 5 a 3 nanómetros, y parte de la mejora acumulada pertenece a TSMC y
  no al diseño. La comparación del Panel A de la Figura~\ref{fig:fom} atenúa
  ese problema pero no lo resuelve.
\item \textbf{La atribución causal de la respuesta de los competidores es
  interpretativa.} Que Intel integrara la memoria en el paquete en 2024 es un
  hecho; que lo hiciera por el M1 es una lectura razonable, no una afirmación
  documentada por la empresa.
\end{enumerate}

\printbibliography

\end{document}
